% ==============================================================================
% CHAPTER 8: CHALLENGES AND FUTURE DIRECTIONS
% ==============================================================================
\section{Discussion}
\label{sec:discussion}

The preceding analysis shows that multi-agent collaborative QA is not limited
by a single algorithmic bottleneck. Its practical effectiveness depends on the
joint interaction of architecture, communication, evaluation, domain
knowledge, human oversight, and deployment constraints. On this basis, this
section first identifies four open challenges and then outlines five future
research directions. The two sets are not intended to correspond one-to-one:
some directions address several challenges simultaneously, whereas others
arise from broader opportunities created by persistent multi-agent
interaction.

\subsection{Open Challenges}
\label{sec:open_challenges}

\textbf{Challenge 1: Architecture adaptation under competing task
requirements and deployment constraints.}
The taxonomy in Section~3 distinguishes flat, hierarchical, and hybrid
architectures according to their different design priorities. Flat
architectures favour decentralised exploration and multi-perspective
verification, hierarchical architectures favour controlled decomposition and
convergence, and hybrid architectures attempt to combine these properties.
Real-world QA tasks, however, are rarely dominated by a single requirement.
They may simultaneously require exploratory discussion, reliable convergence,
low latency, specialised knowledge, and strict risk control. The
architecture-selection stress test in Section~6 illustrates this boundary.
For EduPlanner \citep{zhang2024eduplanner}, open-ended instructional design
favours peer exploration, whereas curriculum-scoring objectives favour a
controlled evaluator--optimizer process. Financial Reflection QA
\citep{fatemi2024financialqa} similarly combines a convergent numerical
objective with iterative multi-critic reflection. In both cases, a
single-dimension decision rule did not reproduce the architecture adopted by
the original system.

This ambiguity extends beyond choosing among three architecture classes. A
deployed service must also decide which agents should participate, how many
are needed, which communication topology to use, how many interaction rounds
to permit, and when to stop. Adding agents may broaden expertise, but it also
increases inference, communication, and tool-use costs and can introduce
redundant or conflicting information. The design problem is therefore
constrained jointly by answer quality, latency, monetary cost, and risk rather
than by professional relevance alone. Static qualitative rules remain useful
for initial screening, but their robustness decreases when several task
requirements compete and no single objective clearly dominates.

\textbf{Challenge 2: Limited observability and process-level evaluation of
multi-agent collaboration.}
Most QA benchmarks emphasise final task outcomes, such as answer accuracy or
completion success \citep{hendrycks2021mmlu,yang2018hotpotqa,
liu2024agentbench,mialon2024gaia}. These measures are necessary, but they
cannot distinguish fundamentally different collaboration processes. A
correct answer may result from two rounds of productive disagreement and
error correction, or from a long sequence of repetitive exchanges followed
by an incidental success. Likewise, a higher aggregate score does not by
itself establish that the participating agents contributed complementary
knowledge rather than restating the same reasoning in different words.

The communication analysis in Sections~3 and~4 further shows why this
distinction matters. In a fully connected broadcast topology, the number of
potential pairwise messages in each round grows quadratically with the number
of agents. More communication can therefore consume substantially more
tokens without guaranteeing a commensurate information gain
\citep{qian2024scaling_multiagent}. Existing task-level scores rarely reveal
which round corrected an error, which role supplied decisive evidence, which
arguments were adopted in the final answer, or whether apparent agreement
reflects genuine convergence or pseudo-consensus. Trace-based analyses such as
MAST demonstrate the value of examining system design, inter-agent
misalignment, and task verification failures \citep{han2024mast}, but such
diagnostics are not yet standard components of collaborative QA evaluation.
Consequently, current evaluations indicate whether a system is correct more
readily than whether its collaboration is efficient, complementary, and
reliable.

\textbf{Challenge 3: Insufficient human control, traceability, and
accountability in high-stakes QA.}
Medical and legal QA require more than a plausible final answer. Evidence
provenance, intermediate reasoning, role-specific contributions, and the
authority to approve or reject consequential actions may all need to be
examined. The medical and legal systems reviewed in Section~5 illustrate the
importance of specialised knowledge and multi-role reasoning
\citep{tang2024medagents,cui2023chatlaw}; they also expose the gap between
benchmark performance and the requirements of accountable deployment. This
gap is particularly important when an incorrect answer can affect a clinical,
legal, or organisational decision.

The HITL architectures discussed in Section~3 provide a foundation for human
participation, but many implementations treat intervention primarily as a
workflow breakpoint, a textual correction, or final approval
\citep{duan2024langgraph_crewai,langgraph2024}. Such mechanisms answer whether
a human can enter the loop, but not when intervention is most effective, what
intermediate state should be editable, or what authority the reviewer should
possess. They may also leave incomplete records of the evidence viewed, the
state changed, and the effect of the intervention on subsequent reasoning.
The central challenge is therefore not merely to include a human in the
workflow, but to provide a degree of control and auditability proportionate to
the risk of the QA task.

\textbf{Challenge 4: Limited adaptability and composability for real-world
deployment.}
Two deployment patterns are likely to coexist in practice. Some applications
require a persistent domain-specific team, such as a multi-specialty medical
QA service that repeatedly handles related cases. Others require an ad hoc
team assembled for a temporary problem, such as urban planning, environmental
assessment, or enterprise decision support. Both patterns require rapid
connection of domain knowledge, tools, and specialised roles. General
frameworks and platforms such as AutoGen, LangGraph, CrewAI, and Dify provide
reusable components for role definition, workflow orchestration, retrieval,
and tool use \citep{wu2024autogen,langgraph2024,crewai2023,dify2024}.
Nevertheless, component-level reuse does not by itself provide task-level
adaptation.

Developers must still specify role responsibilities, knowledge boundaries,
communication protocols, tool permissions, and termination conditions. When
the task or domain changes, many of these decisions must be revisited, while
knowledge sources and evidence formats may not transfer cleanly across
systems. Ad hoc deployments also lack a standard procedure for validating an
assembled team before it is exposed to users. This gap raises the engineering
cost of translating reusable multi-agent components into reliable
task-specific QA services.

These challenges are closely related but operate at different levels.
Architecture adaptation concerns how a system should be configured;
process-level evaluation concerns how collaboration should be measured;
high-stakes deployment concerns how decisions should be supervised and
audited; and composability concerns how a collaborative team can be assembled
for a new task. Addressing them therefore requires a research agenda that
extends beyond isolated improvements to prompts or individual agents.

\subsection{Future Research Directions}
\label{sec:future_directions}

The following directions do not constitute one-to-one responses to the
preceding challenges. Some address several limitations simultaneously, while
continual agent evolution represents a broader opportunity that emerges when
multi-agent QA systems are treated as persistent services rather than
one-off inference pipelines.

\textbf{Direction 1: Process-aware evaluation of collaborative
effectiveness.}
A first priority is to extend evaluation from the final answer to the complete
collaboration trace. Each interaction could be represented as a structured
event containing the speaker, claim, supporting evidence, confidence,
response target, and whether the information was retained in the final
answer. This representation would support metrics such as the marginal gain
in answer quality after each round, the adoption rate of role-specific
arguments, the error-correction rate relative to a single-agent baseline, the
proportion of repeated or non-contributory messages, and the tokens, time, and
tool calls required to reach a fixed quality level. Trace taxonomies such as
MAST offer a starting point for annotating failures \citep{han2024mast}, but
collaborative QA also requires positive measures of useful disagreement,
complementarity, and convergence.

Validation should compare single-agent, fixed multi-agent, and dynamically
configured systems under the same base model, knowledge source, and inference
budget. Human assessment can be used on a subset of traces to determine
whether an apparent improvement reflects genuinely new evidence or merely
linguistic reformulation. The resulting process metrics should then be tested
for their ability to predict final reliability and failure risk across
multi-hop, open-domain, and domain-specific QA tasks. Such a benchmark would
provide a common measurement layer for the remaining research directions.

\textbf{Direction 2: Multi-objective orchestration of agent selection,
interaction, and termination.}
Architecture selection can be reformulated as a dynamic orchestration problem
rather than a one-time classification decision. A controller could estimate
task complexity, domain requirements, risk level, current disagreement, and
evidence sufficiency, and then decide which agents to invoke, how they should
communicate, whether another specialist is needed, and when collaboration
should terminate. For example, a high-confidence agreement supported by
independent evidence could trigger early stopping, whereas unresolved
disagreement or weak evidence could activate an additional critic or domain
expert. Adaptive systems such as MDAgents provide an initial example of
changing collaboration according to problem difficulty
\citep{kim2024mdagents}, while broader scaling results indicate that the
benefit of adding agents depends on both task and coordination design
\citep{qian2024scaling_multiagent}.

The controller should optimise answer quality jointly with latency, token
consumption, monetary cost, and tool-use overhead. Evaluation could compare
static, rule-based, and learned orchestration under multiple resource budgets,
reporting both absolute performance and the Pareto frontier between quality
and resource use. This design would reveal whether dynamic orchestration can
preserve answer quality while reducing unnecessary participants and
interaction rounds. It would also convert the qualitative selection guide in
Section~6 into an empirically testable decision policy.

\textbf{Direction 3: Continual co-evolution of agent knowledge and group
collaboration strategies.}
Persistent multi-agent QA systems should not necessarily restart every case
with the same knowledge and collaboration policy. At the individual level,
an agent could use verified cases, retrieved updates, and human corrections
to refine its domain memory, record recurring error types, and calibrate the
boundary of its expertise. Existing work on agent memory and reflective
feedback provides useful components for such updating
\citep{zhang2024memory_survey,shinn2023reflexion}. To prevent unverified
content from becoming persistent knowledge, each update should retain its
source, timestamp, confidence, and validation status.

At the group level, the system could learn which agents are reliable for
particular question types, which combinations produce complementary evidence,
and which interaction patterns tend to repeat errors. These histories could
inform dynamic speaking order, routing, voting weights, and stopping rules.
For example, a medical system that repeatedly observes reliable pathology
judgements from one specialist agent could assign greater weight to that
agent on related cases while avoiding unnecessary debate. This extends
memory-based individual adaptation towards learned team coordination, an
opportunity suggested by persistent-agent environments
\citep{park2023generative_agents,chen2024agentverse}.

Longitudinal evaluation is essential because short benchmark runs cannot
measure evolution. Experiments should present sequential tasks with changing
knowledge and monitor adaptation speed, knowledge retention, long-term
accuracy, and collaboration cost. They should also test safeguards against
error propagation, collective bias reinforcement, catastrophic forgetting,
privacy leakage, and unstable shifts in agent authority. Comparing individual
updating, group-policy updating, and their joint use would show whether the
two forms of evolution are complementary.

\textbf{Direction 4: Declarative and low-code composition of task-specific
multi-agent QA systems.}
To reduce deployment effort, future platforms could allow users to declare a
task objective, domain knowledge sources, role responsibilities, tool
permissions, risk level, and resource budget. The platform could then generate
an initial team configuration containing agent roles, communication topology,
knowledge interfaces, discussion protocol, and termination rules. Reusable
templates could support both persistent domain teams and temporary
problem-oriented teams. Current orchestration frameworks already provide
several of the required components \citep{wu2024autogen,langgraph2024,
crewai2023,dify2024}; the research challenge is to compose and validate them
at the task level rather than expose them only as separate development
features.

Domain knowledge should also enter collaboration through structured
interfaces rather than prompt text alone. Retrieval adapters, knowledge
graphs, typed evidence objects, confidence estimates, and provenance metadata
could allow agents to exchange claims together with their evidential status.
This approach is more interoperable than assuming direct access to
model-specific hidden representations and is consistent with the structured
communication and memory mechanisms reviewed in Section~4. Before deployment,
an automatically assembled team could be tested on historical cases or
sandbox tasks to identify missing expertise, unsafe tool permissions, and
non-convergent interaction patterns.

Evaluation should measure the time and code required to configure a new
system, performance loss during cross-domain transfer, and the proportion of
configuration errors detected before deployment. Comparisons among manual
engineering, template-based assembly, and automatically generated
configurations would establish whether low-code composition reduces effort
without sacrificing reliability.

\textbf{Direction 5: Risk-adaptive human oversight and auditable
intervention.}
Human oversight should vary with task risk, evidence uncertainty, and
inter-agent disagreement rather than being inserted at a fixed point in every
workflow. A system could trigger review before task decomposition, after
evidence retrieval, when agents reach a material disagreement, before a
high-impact tool action, or before the final answer is released. In
high-risk cases, reviewers should be able to inspect and revise intermediate
states rather than only accept or reject the final output. Each intervention
should generate an audit record containing the reviewer, evidence inspected,
state changed, reason for the change, and effect on subsequent reasoning.

This design can be evaluated by constructing QA tasks at several risk levels
and comparing fixed breakpoints, uncertainty-triggered review,
disagreement-triggered review, and continuous human supervision. Relevant
outcomes include error-interception rate, decision time, human workload,
answer quality, user trust, and the completeness of responsibility tracing.
Crucially, experiments should test whether additional transparency improves
error detection rather than merely increasing perceived trust. Such evidence
would provide risk-specific engineering guidance for medical, legal, and
other regulated applications.

\begin{table*}[pos=t]
  \centering
  \caption{Operational research agenda for multi-agent collaborative QA.}
  \label{tab:research_agenda}
  \small
  \begin{tabularx}{\textwidth}{p{2.8cm} X p{3.5cm} p{3.1cm}}
    \toprule
    \textbf{Research direction} &
    \textbf{Operational mechanism} &
    \textbf{Key evaluation criteria} &
    \textbf{Validation setting} \\
    \midrule
    Process-aware evaluation &
    Structured interaction traces and contribution analysis &
    Marginal gain, correction rate, redundancy, and resource cost &
    Multi-hop and debate-based QA \\[2pt]
    Multi-objective orchestration &
    Dynamic agent selection, routing, escalation, and termination &
    Quality, latency, token cost, and Pareto efficiency &
    Budget-constrained QA \\[2pt]
    Continual co-evolution &
    Verified knowledge updating and learned reliability profiles &
    Adaptation speed, retention, long-term accuracy, and stability &
    Sequential domain tasks \\[2pt]
    Declarative composition &
    Automatic role, tool, knowledge, and protocol configuration &
    Setup time, code reduction, transfer performance, and pre-deployment errors &
    Cross-domain deployment \\[2pt]
    Risk-adaptive oversight &
    Uncertainty- and disagreement-triggered human intervention &
    Error interception, workload, trust, and traceability &
    Medical, legal, and other regulated QA \\
    \bottomrule
  \end{tabularx}
\end{table*}

\subsection{Limitations of This Survey}
\label{sec:survey_limitations}

This survey has four principal limitations. First, the literature search
extends to mid-2026, and the speed of development means that later or
insufficiently indexed systems cannot be represented. Second, architecture
classification depends on the level of implementation detail reported by the
original authors. Systems that dynamically switch topology or combine several
coordination mechanisms may therefore remain ambiguous at category
boundaries. Third, the ten-case comparison in Section~6 is a small diagnostic
stress test rather than independent confirmation of general validity; broader
testing across tasks, models, and resource budgets is required. Fourth, the
architecture-selection guide remains qualitative. Its decision factors have
not yet been operationalised as a weighted model learned or calibrated on a
large labelled system corpus. These limitations bound the claims that can be
made for the proposed taxonomy and selection guide, while motivating their
future empirical refinement.
